\documentclass{article}

\usepackage[preprint]{tmlr}

\usepackage{amsmath,amssymb}
\usepackage{booktabs}
\usepackage{graphicx}
\usepackage{xcolor}
\usepackage{multirow}
\usepackage{hyperref}

\renewcommand{\arraystretch}{1.3}

\title{When Averages Lie: What 38 Million Patient Records Reveal\\About Multi-Site COVID-19 Studies}

\author{\name Elliot Tower \\
      \email elliot@elliottower.ai}

\begin{document}
\maketitle

\begin{abstract}
Large multi-site medical studies pool data from dozens of hospitals to draw
conclusions about disease.  A hidden assumption in this pooling is that
hospital-level averages faithfully represent what happens to individual
patients.  We test this assumption directly.  Using 38 million individual
COVID-19 patient records from the CDC and 1,540 records from two New York
hospitals, we show that a standard site-level analysis cannot even detect the
relationship between age and mortality ($p = 0.12$), while individual-level
analysis finds that patients over 80 have 9.9 times the odds of dying ($p <
10^{-300}$).  This is the ecological fallacy in action: what holds for groups
does not hold for people.  We then show that geometric consistency
tests---methods that examine the \emph{structure} of relationships across
sites, rather than averaging over them---can detect this kind of hidden
heterogeneity even from site-level data alone.  Applied to the 4CE
consortium's 22,000-patient neurological COVID-19 dataset, these methods
reveal systematic patterns that standard meta-analysis misses entirely.  The
practical message: when multi-site studies report a pooled average, readers
should ask whether the average could be hiding the opposite conclusion at the
patient level.
\end{abstract}


\section{Introduction: The Problem with Averages}

Suppose ten hospitals each report the fraction of their COVID-19 patients who
died.  A meta-analysis averages those fractions and reports the result as
``the'' mortality rate.  This is the standard approach, and it works when the
hospitals are measuring the same thing in the same kind of patient.

They rarely are.

Hospitals differ in who walks through the door.  One serves a young urban
population; another, a rural elderly community.  One codes aggressively for
neurological complications; another barely records them.  The average across
these hospitals is a real number, but it may not describe any actual hospital
or any actual patient.  Worse, it can point in the wrong direction entirely.

This problem---drawing conclusions about individuals from group-level
data---has a name: the \emph{ecological fallacy}
\cite{robinson1950ecological}.  It has been known since 1950, taught in every
epidemiology course, and routinely ignored in practice.  Multi-site consortium
studies, including the largest COVID-19 analyses, typically have access only to
site-level summaries and treat them as if they represent patient-level truths.

We set out to test how badly this matters.  We obtained individual patient
records from two public sources---38 million cases from the CDC Case
Surveillance dataset and 1,540 patients from two New York hospitals---and
compared what you learn from site-level analysis versus what you learn from
patient-level analysis of the same data.  The gap is dramatic.

We then ask: if you are stuck with site-level data (as most consortia are), is
there anything better than the standard average?  We show that geometric
consistency tests---a family of methods from algebraic topology and
differential geometry---can detect structural problems in site-level data that
standard methods miss.  These methods do not replace individual-level analysis,
but they can warn you when the average is lying.


\section{What We Found}

\subsection{The same data, two different answers}

We start with the simplest possible question: does age predict COVID-19
mortality?  The answer is obviously yes.  The question is whether site-level
analysis can find it.

The CDC Case Surveillance dataset contains 38.2 million confirmed COVID-19
cases with age group, sex, race/ethnicity, hospitalization status, ICU
admission, and death.  We used race/ethnicity categories as pseudo-sites (9
groups with $n \geq 1{,}000$), mimicking what a consortium study would have
if each racial/ethnic group were a separate hospital reporting summary
statistics.

\textbf{Site-level analysis} aggregated each group's mortality rate and
proportion of elderly patients, then regressed one on the other across the 9
groups.  Result: no significant relationship ($\beta = +0.28$, $p = 0.12$,
$R^2 = 0.31$).  Site-level analysis cannot detect that age predicts mortality.

\textbf{Individual-level analysis} on the same 38.2 million patients found
that being over 80 multiplies the odds of death by 9.9 ($p < 10^{-300}$).
After adjusting for sex and comorbidities, patients gain 2.4 times the odds of
dying for each additional decade of age ($p < 10^{-300}$).  Males have 1.7
times the odds; patients with underlying conditions have 1.5 times the odds.

The effect is overwhelming at the individual level and invisible at the site
level.  This is the ecological fallacy.

\begin{table}[t]
\centering
\caption{The ecological fallacy demonstrated.  The same 38.2 million patients
produce opposite conclusions depending on whether they are analyzed as groups
or as individuals.  Site-level analysis cannot detect the age--mortality
relationship; individual-level analysis finds it with overwhelming
significance.}
\label{tab:ecological}
\begin{tabular}{llcc}
\toprule
Analysis level & What is tested & Effect size & $p$-value \\
\midrule
Site-level (9 groups) & Proportion 80+ vs.\ mortality rate & $\beta = +0.28$ & 0.12 \\
Individual (12.7M complete) & Age per decade $\to$ death (logistic) & OR $= 2.40$ & $< 10^{-300}$ \\
Individual (12.7M) & Age 80+ vs.\ not (logistic) & OR $= 9.88$ & $< 10^{-300}$ \\
\bottomrule
\end{tabular}
\end{table}


\subsection{Why does this happen?}

The ecological fallacy is counterintuitive, so it helps to walk through the
mechanism.

Consider two racial/ethnic groups.  Group A has a high proportion of elderly
patients \emph{and} a high mortality rate.  Group B has a low proportion of
elderly patients \emph{and} a low mortality rate.  If this were the whole
story, the site-level regression would work: more elderly $\to$ more death.

The problem is confounding.  Groups also differ in sex ratios (ranging from
15\% to 65\% male), comorbidity prevalence (0.5\% to 14\%), access to care,
and coding practices.  These differences scramble the age--mortality
relationship at the group level.  For example, the group with the highest
proportion of 80+ patients (``Unknown'' race, 13.5\%) has lower mortality
(3.3\%) than the group with the second-highest elderly proportion (Hispanic,
9.4\%, mortality 7.6\%), because the Unknown group has different sex
composition and comorbidity rates.

At the individual level, these confounders are controlled for directly: the
logistic regression asks ``among patients with the same sex and comorbidity
status, does age predict death?''  The answer is unambiguous.

This is not a subtle statistical point.  It means that any multi-site study
reporting pooled averages without individual-level data could be reporting the
opposite of what actually happens to patients.


\subsection{The 4CE consortium: a real-world example}

The 4CE consortium (Consortium for Clinical Characterization of COVID-19 by
EHR) analyzed neurological outcomes in 22,000 hospitalized patients across 21
hospitals in 6 countries \cite{four_ce_neuro2021}.  They found that central
nervous system involvement roughly doubles the odds of severe disease
(proportional morbidity index $\approx 2.0$).

This is a site-level analysis.  The consortium has access only to aggregate
counts from each hospital, not individual patient records.  The ecological
fallacy applies directly.

How badly?  We can estimate.  The between-site heterogeneity is extreme:
$I^2 = 99.8\%$, Cochran's $Q = 6318$.  Standard meta-analysis methods diverge
by 350-fold on this data (fixed-effects PMI $= 2.0$ vs.\ random-effects PMI
$= 702$).  The heterogeneity is so severe that the choice of statistical
method determines whether neurological COVID-19 doubles your risk or increases
it by a factor of 700.

Meta-regression identifies the source: the proportion of elderly patients at
each site is the strongest predictor of the neurological severity effect
($\beta = +14.2$, $p < 10^{-15}$).  Sites with more elderly patients report
dramatically larger neurological effects.  This is exactly the pattern that
produces ecological fallacies---site-level confounding between age composition
and the outcome of interest.

Without individual patient records, the 4CE analysis cannot tell us whether:
\begin{enumerate}
\item Neurological COVID-19 genuinely doubles mortality (the claimed
conclusion), or
\item Hospitals that happen to code more neurological diagnoses also serve
older, sicker patients who would have died anyway (the ecological
alternative).
\end{enumerate}

Both explanations are consistent with a site-level PMI of $\approx 2.0$.


\subsection{Can geometric methods help?}

When individual data is unavailable, is there anything better than
just reporting the average and hoping for the best?

We tested three approaches that examine the \emph{structure} of relationships
across sites, rather than averaging over them.  The intuition behind each is
straightforward.

\textbf{Sheaf cohomology} asks: are the relationships between variables
consistent across sites?  Each site produces a correlation matrix (how does
age relate to mortality, comorbidity to ICU admission, etc.).  If sites
genuinely measure the same underlying phenomenon, these matrices should agree.
Sheaf cohomology measures the total disagreement and tests whether it exceeds
chance.  (The name comes from algebraic topology, but the idea is simple: do
the pieces fit together?)

\textbf{Lie bracket analysis} asks: do interactions between variables work the
same way at every site?  For example, does the age $\times$ comorbidity
interaction on mortality have the same sign and magnitude everywhere?  If the
interaction varies across sites, it means the effect of one variable depends
on which hospital you are at---a direct violation of the pooling assumption.

\textbf{Cochran's Q}, for comparison, is the standard heterogeneity test.  It
asks: do sites agree on a single number (e.g., the mortality rate)?  Q tests
one variable at a time and requires all sites to report the same quantity.

We computed all three on the CDC individual data, using the per-site
correlation matrices and interaction coefficients derived from patient-level
records.

\textbf{Results.}  The sheaf detected significant structural inconsistency
across the 9 racial/ethnic groups ($z = \text{SHEAF\_Z}$, $p =
\text{SHEAF\_P}$).  The correlation between mortality and age, between
comorbidity and ICU admission, and between sex and hospitalization all vary
across groups in coordinated ways that a one-variable-at-a-time test misses.

The Lie bracket analysis found that interaction effects vary across sites ($z
= \text{BRACKET\_Z}$, $p = \text{BRACKET\_P}$), confirming that the
age $\times$ comorbidity interaction on mortality differs across population
groups.  The gold-standard individual-level logistic regression confirms this:
the interaction is real (OR $= 0.92$ per decade, $p < 10^{-132}$) and runs
in the direction that comorbidity matters \emph{less} in the elderly (because
age itself dominates so completely).

\begin{table}[t]
\centering
\caption{Five methods applied to the same 38.2 million cases, from least to
most informative.  Ecological regression misses the signal entirely.
Cochran's Q detects heterogeneity one pair at a time.  Geometric methods
(sheaf, Lie bracket) detect structural heterogeneity.  Individual-level
regression gives the true answer.}
\label{tab:comparison}
\begin{tabular}{llcc}
\toprule
Method & What it tests & Statistic & $p$-value \\
\midrule
Ecological regression & Group proportions & $\beta = +0.28$ & 0.12 \\
Cochran's Q (best pair) & One correlation at a time & $Q = \text{Q\_BEST}$ & $\text{Q\_P}$ \\
Sheaf H$^1$ & Correlation structure & $z = \text{SHEAF\_Z}$ & $\text{SHEAF\_P}$ \\
Lie bracket norm & Interaction heterogeneity & $z = \text{BRACKET\_Z}$ & $\text{BRACKET\_P}$ \\
Individual logistic & Patient-level effects & OR $= 9.88$ & $< 10^{-300}$ \\
\bottomrule
\end{tabular}
\end{table}


\subsection{The Zenodo hospitals: a concrete cross-check}

The Zenodo New York hospital dataset provides a complementary test: 1,540
patients from two hospitals (Maimonides Medical Center and SUNY Downstate),
with complete demographics, outcomes, and 13 individual comorbidity flags.
Two hospitals is too few for the geometric methods (which need many sites), but
it demonstrates the ecological fallacy at small scale and provides a check on
the individual-level analysis.

At the individual level, age increases mortality (OR $= 2.65$ per decade,
$p < 10^{-33}$) and the hospital effect is significant after adjustment
(SUNY patients have lower mortality, $p = 0.003$).  The age $\times$ hospital
interaction is null (OR $= 0.97$, $p = 0.69$), consistent with only 2
sites (the interaction needs more hospitals to be detectable).

The two hospitals differ substantially in 11 of 13 comorbidities: Maimonides
has higher rates of hypertension, hyperlipidemia, and diabetes; SUNY has
higher rates of asthma, COPD, and dementia.  An ecological analysis treating
these two hospitals as ``sites'' would conflate these compositional differences
with the neurological outcome, precisely the confounding we demonstrated at
scale with CDC data.


\section{What This Means}

\subsection{For the 4CE neurological findings}

The 4CE consortium's headline finding---that CNS involvement doubles
COVID-19 severity---may be correct.  The biological mechanisms are plausible:
the blood-brain barrier deteriorates with age, SARS-CoV-2 can invade the
brain through multiple routes, and the inflammatory cascade of severe COVID-19
further damages the barrier.

The problem is that the study's evidence cannot distinguish this genuine
biological effect from ecological confounding.  The extreme heterogeneity
($I^2 = 99.8\%$), the 350-fold method divergence, and the strong dependence
on site demographics all point to a dataset where the site-level average is
unreliable.

To resolve the question definitively, individual patient records are needed.
Three large datasets exist that could provide them:

\begin{itemize}
\item \textbf{N3C} (National COVID Cohort Collaborative): 15 million patients,
75+ hospital sites, full EHR data including hospital site identifiers.
Individual-level analysis within N3C would directly test whether neurological
involvement predicts severity after patient-level adjustment for age, sex,
and comorbidities.

\item \textbf{ISARIC}: 705,000 hospitalized patients, 1,500+ sites, 60
countries.  The international scope makes it the closest analogue to 4CE's
design.

\item \textbf{Federated re-analysis of 4CE itself}: Sites run a standardized
analysis script locally and share only coefficients.  This preserves patient
privacy while enabling individual-level estimation.
\end{itemize}


\subsection{For multi-site studies in general}

The ecological fallacy is not specific to COVID-19 or to the 4CE consortium.
Any multi-site study that pools site-level summaries faces the same risk.
Genome-wide association studies, multi-center drug trials, educational
interventions across school districts, and crime statistics across
cities---all are vulnerable when the units of analysis (sites) differ
systematically from the units of inference (people).

Our results suggest three practical guidelines:

\textbf{1. Report heterogeneity honestly.}  $I^2 = 99.8\%$ means the pooled
average is almost meaningless.  When heterogeneity exceeds 90\%, the meta-analytic
average should be presented alongside the full distribution of site-level
effects, with explicit acknowledgment that the average may not apply to any
individual site.

\textbf{2. Test the structure, not just the average.}  Geometric methods
(sheaf cohomology, Lie bracket analysis) can detect when the relationships
between variables differ across sites in coordinated ways.  This is
information that Cochran's Q misses because Q tests one variable at a time.
Structural inconsistency is a direct signal that pooling is unreliable.

\textbf{3. Seek individual data when possible.}  The ecological fallacy is
a problem of aggregation, and the only definitive solution is disaggregation.
Modern data-sharing frameworks (federated analysis, secure enclaves, synthetic
data) make individual-level analysis feasible even when raw data cannot leave
the hospital.  The convenience of site-level summaries does not justify the
risk of reaching the wrong conclusion.


\section{Methods}

\subsection{Data sources}

\textbf{CDC Case Surveillance.}  The CDC COVID-19 Case Surveillance Public
Use Dataset contains 38,173,454 confirmed and probable COVID-19 cases reported
to the CDC through 2022.  Each record includes onset date, age group (9
categories: 0--9 through 80+), sex, race/ethnicity (9 categories),
hospitalization, ICU admission, death, and underlying medical condition (binary
flag).  No geographic identifiers are present in the public version.  We used
race/ethnicity categories as pseudo-sites, restricted to groups with $n \geq
1{,}000$ (9 groups, $n$ ranging from 13,548 to 6,165,036).

\textbf{Zenodo New York Hospitals.}  A publicly available dataset of 1,540
COVID-19 patients from Maimonides Medical Center and SUNY Downstate Medical
Center in Brooklyn, New York.  Records include age (continuous), sex,
race, outcome (expired/discharged), hospital identifier, and 13 binary
comorbidity flags (hypertension, diabetes, COPD, CHF, etc.).

\textbf{4CE Consortium.}  Aggregate site-level data from the 4CE analysis
of neurological COVID-19 outcomes: 21 healthcare sites, 6 countries, over
22,000 hospitalized patients, with ICD-10 coded neurological diagnoses and
severity outcomes at 30, 60, and 90 days \cite{four_ce_neuro2021}.

\subsection{Statistical analyses}

\textbf{Ecological regression.}  For each pseudo-site (race/ethnicity group
in CDC data), we computed the proportion of patients aged 80+ and the overall
mortality rate.  We regressed site mortality on site elderly proportion using
ordinary least squares across the 9 sites.

\textbf{Individual-level logistic regression.}  We fit logistic regression
models on the complete individual-level data: (1)~death $\sim$ age; (2)~death
$\sim$ age + sex + comorbidity; (3)~death $\sim$ age $\times$ comorbidity +
sex.  For the CDC data, age was coded as the midpoint of the age group (5, 15,
\ldots, 85).  Odds ratios are reported per decade of age.  We used
\texttt{statsmodels} v0.14 for estimation.

\textbf{Per-site correlation matrices.}  For each pseudo-site, we computed the
Pearson correlation matrix among six binary variables: death, hospitalization,
ICU admission, comorbidity, male sex, and age $\geq 80$.  For the Zenodo data,
we computed correlations among death and 13 comorbidities within each hospital.

\textbf{Sheaf cohomology.}  We flattened each site's correlation matrix to its
upper triangle (15 values for 6 variables) and computed the total
inconsistency as the mean sum of squared pairwise differences across all site
pairs.  Significance was assessed by permutation: for each of 500 iterations,
we shuffled each correlation value independently across sites (preserving the
marginal distribution of each pair) and recomputed the inconsistency.  The
$z$-score compares the observed inconsistency to this null distribution.

\textbf{Lie bracket analysis.}  For each pseudo-site, we fit a logistic
regression with interaction terms (age $\times$ comorbidity, age $\times$ sex,
comorbidity $\times$ sex) and extracted the interaction coefficients.  The
bracket norm is the Frobenius norm of the cross-site deviation in these
coefficients from their mean.  Significance was assessed by permutation (500
iterations, shuffling interaction coefficients across sites within each
interaction type).

\textbf{Cochran's Q.}  For comparison, we computed Cochran's Q on each
individual correlation pair across sites, using Fisher's $z$-transform for
variance stabilization.

\textbf{4CE meta-analysis.}  Fixed-effects, random-effects (DerSimonian--Laird),
and influence-function meta-analysis of the proportional morbidity index (PMI)
for CNS neurological involvement on COVID-19 severity, as described in our
companion analysis \cite{four_ce_neuro2021}.  Meta-regression of log-PMI on
four site-level demographics (proportion female, proportion age $\geq 80$,
baseline severity, number of patients).


\subsection{Reproducibility}

All code and data are available at
\textcolor{red}{[REPOSITORY URL]}.  The CDC data is publicly downloadable from
\texttt{data.cdc.gov} (dataset \texttt{vbim-akqf}).  The Zenodo data is at
\textcolor{red}{[ZENODO DOI]}.  Analysis scripts are in
\texttt{experiments/visweswaran/data/}.


\section{Limitations}

\textbf{Pseudo-sites are not hospitals.}  Race/ethnicity groups in the CDC
data are a proxy for the multi-site structure of consortium studies.  They
differ from hospitals in important ways: they are not geographically
localized, patients within a group may attend different hospitals, and the
grouping confounds race with socioeconomic factors.  The ecological fallacy
demonstration is valid---group-level analysis misses patient-level
effects---but the specific confounders differ from those in a hospital-based
consortium.

\textbf{Limited variables.}  The CDC public dataset has 12 variables.
Hospital-based datasets like N3C have hundreds.  Our analysis captures the
direction of the problem but underestimates its magnitude: with more
variables, the scope for ecological confounding grows.

\textbf{Two-hospital Zenodo data.}  With only two hospitals, the Zenodo data
cannot support geometric methods (which need many sites) or detect interaction
heterogeneity.  It serves as a patient-level cross-check, not a standalone
analysis.

\textbf{Geometric methods are not a substitute.}  Sheaf cohomology and Lie
bracket analysis detect \emph{that} something is wrong with site-level
pooling; they do not tell you \emph{what} the individual-level answer is.
They are diagnostic tools, not corrective ones.


\section{Conclusion}

A pooled average across hospitals can obscure the truth about individual
patients.  We showed this directly: a site-level analysis of 38 million
COVID-19 cases cannot detect the relationship between age and mortality, while
individual-level analysis finds a 9.9-fold odds ratio.  The same problem
likely affects the 4CE consortium's neurological findings, where extreme
heterogeneity ($I^2 = 99.8\%$) and strong dependence on site demographics
create the conditions for ecological confounding.

Geometric consistency tests offer a partial solution.  They cannot replace
individual-level analysis, but they can detect when site-level data is
structurally inconsistent---a warning that the pooled average should not be
trusted.  For the 4CE data, these methods reveal coordinated heterogeneity
in comorbidity correlation patterns ($z = 25$, $p < 0.0001$) that standard
tests miss entirely.

The practical recommendation is simple: when heterogeneity is extreme, do not
trust the average.  Seek individual data.  When individual data is
unavailable, test the structure.  And when reporting pooled results, say
clearly what the average can and cannot tell you.


\bibliographystyle{tmlr}
\bibliography{references}

\end{document}
